\documentclass[12pt]{report}

\usepackage{geometry}
\usepackage{fontspec}
\usepackage{amsmath}
\usepackage{fancyhdr}
\usepackage{lastpage}
\usepackage[hidelinks]{hyperref}

\setmainfont{Palatino Linotype}
\newfontfamily{\telcontar}[Renderer=Graphite]{Tengwar Telcontar}[Scale=2]
\newfontfamily{\ttbig}[Renderer=Graphite]{Tengwar Telcontar}[Scale=3]
\newfontfamily{\ds}{Doulos SIL}
\setlength{\parskip}{1em}
\pagestyle{fancy}
\fancyhf{}
\lfoot{Page \thepage \hspace{1pt} of \pageref{LastPage}}
\renewcommand{\headrulewidth}{0pt}
\renewcommand{\thesection}{\arabic{section}}
\renewcommand{\thefootnote}{\arabic{footnote}}
\geometry{left=2.5cm,right=2.5cm,top=2.5cm,bottom=3.5cm}

\begin{document}
\title{\font\h=cmr12 at 36pt \h Phi\\
\Large A constructed language}
\author{Daniel Ellison}
\date{\today}
\maketitle
\section{Introduction}
\textsc{Phi} is a constructed language. A constructed language is one whose phonology, grammar, and vocabulary, instead of having developed naturally, are consciously devised\footnote{\url{https://en.wikipedia.org/wiki/Constructed\_language}}. The sound of the name \emph{Phi} embodies the simplicity and softness of the language itself. \emph{Phi} (\ds ɸ) is also beauty in mathematics. It is

\begin{quote}
  \emph{\kern-0.5em “the golden ratio, an irrational number with a value of approximately 1.618033988 which expresses the relationship that the sum of two quantities is to the larger quantity as the larger is to the smaller.”} - Wiktionary \footnote{\url{https://en.wiktionary.org/wiki/\%CF\%86}}
\end{quote}

It is thought that the value of \emph{phi} is one of the fundamental characteristics of nature; in the language Phi, nature is fundamental to its being. Listening to spoken Phi one is sometimes reminded of the sounds of the natural world. There is also a musicality and rhythm to sentences in the language.
\clearpage
\section{Phonology}
\begin{table}[h]
  \caption{\label{consonants}Consonants}
  \begin{tabular}
    {|c|c|c|c|c|c|c|c|}
    \hline
    & Bilabial & Labial & Dental & Alveolar & Postalveolar & Velar & Glottal \\
    \hline
    Plosive & p & & & t & & & \\
    \hline
    Nasal & m & & & n & & & \\
    \hline
    Trill & & & & r & & & \\
    \hline
    Tap or flap & & & & ɾ & & & \\
    \hline
    Fricative & \ds ɸ & \ds ʍ & \ds θ & s & \ds ʃ & \ds ʍ & h \\
    \hline
    Approximate & & w & & & & w & \\
    \hline
    Lateral\\approximant & & & & l & & & \\
    \hline
  \end{tabular}
\end{table}

\noindent
\begin{table}[h]
  \caption{\label{vowels}Vowels}
  \begin{tabular}
    {|c|c|c|c|}
    \hline
    & Front & Central & Back \\
    \hline
    Close & i & & \\
    \hline
    Mid & e & & o \\
    \hline
    Open & & a & \\
    \hline
  \end{tabular}
\end{table}
\clearpage
\section{The Tengwar}
Phi can be written in either of two ways: using the standard Latin alphabet or the Tengwar script. Tengwar was invented by J. R. R. Tolkien and first described in Appendix E of the third book in his trilogy, \emph{The Lord of the Rings}. There are numerous sources of information available concerning the Tengwar so details and history of the script will not be explored in this document.

The look and feel of the Tengwar characters (\emph{tengwa}) and its diacritical marks (\emph{tehta}) match the ideals and philosophies behind this language. Phi has a small set of consonants and vowels so it is relatively straightforward to map its graphemes to Tengwar characters. The Unicode font \emph{Tengwar Telcontar}\footnote{\url{http://freetengwar.sourceforge.net/tengtelc.html}} is used in this manual to render Phi text in Tengwar script. A custom keyboard layout\footnote{\url{http://freetengwar.sourceforge.net/keylayouts.html}} was employed to facilitate typing the Tengwar characters.

Tables 3 and 4 outline the mapping between the graphemes used in Phi and the corresponding Tengwar characters, as well as the keypresses necessary to produce the characters using the font and the keyboard layout mentioned previously. An application or utility that supports Graphite fonts\footnote{\url{https://scripts.sil.org/cms/scripts/page.php?site\_id=projects\&item\_id=graphite\_fonts}} is required for the font to render correctly. Experiments indicate that LibreOffice Writer v7.1.2.2 renders Tengwar Telcontar properly. A link to details of this requirement can be found in the first footnote below.

Written Phi uses the Tengwar’s Quenya mode. Vowel diacriticals are placed on the consonant that precedes the vowel. Since all Phi words start with a consonant, there is never a need to use an initial vowel carrier. Words in Phi often contain two consecutive vowels, but a vowel carrier is still unnecessary: the second vowel is placed below the same consonant. The letters would then be read as consonant, upper vowel, then lower vowel. This and the use of a single Tengwar character for consonant doubles such as ph or sh cause Phi script to be very compact and efficient.
\clearpage
The Phi word for “love” is \emph{lothea}. In the Tengwar script it looks like this:

\bigskip
\leftskip=0.66in
\ttbig  \normalfont
\bigskip

\leftskip=0in
This six-letter word in latin characters requires only two \emph{tengwa} to represent it. Many words in Phi are two to three letters long. This occurs mainly in parts of speech such as particles, prepositions, or pronouns. These words are almost always a single \emph{tengwa} with diacriticals in Tengwar script. All of this results in a very compact but clear Phi writing system.

This is the phrase “I love you” in its simplest vernacular form in Phi:

\bigskip
\leftskip=0.66in
\ttbig  \normalfont
\bigskip

\leftskip=0in
That is a beautiful sentiment rendered in a beautiful script. In more formal communication in Phi, particles are used to mark parts of speech in order to reduce ambiguity. As well, word markers are often used to distinguish multisyllabic words for clarity. The Phi Tengwar punctuation is outlined in Table 5. As an example, the phrase “I love you” would be written formally as:

\bigskip
\leftskip=0.66in
\ttbig ⸱⸱⸱⸱⸱ \normalfont
\clearpage
\leftskip=0in
\section{Examples}
Following are several examples of Phi phrases in formal mode. They are given first in English, then in romanized text, and finally in the Tengwar script.
\setlength{\parskip}{1.5em}

{\parindent0pt
My name is Daniel. -- \emph{si phiathe soa mia na thanie te phi}
\setlength{\parskip}{0.3em}

\telcontar ⸱⸱⸱⁘⸱⸱⸱ \normalfont
\setlength{\parskip}{1.5em}

I love you. -- \emph{mia na thi te lothea}
\setlength{\parskip}{0.3em}

\telcontar ⸱⸱⸱⸱ \normalfont
\setlength{\parskip}{1.5em}

The sky is blue. -- \emph{whetho na mipho te phi}
\setlength{\parskip}{0.3em}

\telcontar ⸱⸱⁘⸱ \normalfont
\setlength{\parskip}{1.5em}

They like food. -- \emph{lo sha na noshea te lephe}
\setlength{\parskip}{0.3em}

\telcontar ⸱⸱⸱⸱⸱ \normalfont
\setlength{\parskip}{1.5em}

I like it. -- \emph{mia na sha te lephe}
\setlength{\parskip}{0.3em}

\telcontar ⸱⸱⸱⸱ \normalfont
\setlength{\parskip}{1.5em}

I do not like green! -- \emph{ho mia na hashe te me lephe}
\setlength{\parskip}{0.3em}

\telcontar ⸱⸱⸱⸱⸱⸱ \normalfont
\setlength{\parskip}{1.5em}

I liked my friend. -- \emph{mia na miwho soa mia te li lephe}
\setlength{\parskip}{0.3em}

\telcontar ⸱⸱⸱⸱⸱⸱⸱ \normalfont
\setlength{\parskip}{1.5em}

The food is ours. -- \emph{noshea na soa lo mia te phi}
\setlength{\parskip}{0.3em}

\telcontar ⸱⸱⸱⸱⸱⸱ \normalfont
%\setlength{\parskip}{1.5em}

\clearpage
Our food is good. -- \emph{noshea soa lo mia na shishoa te phi}
\setlength{\parskip}{0.3em}

\telcontar ⸱⸱⸱⸱⸱⸱⸱ \normalfont
\setlength{\parskip}{1.5em}

Are you in good health? -- \emph{wha thi na ropho shishoa te la phe}
\setlength{\parskip}{0.3em}

\telcontar ⸱⸱⸱⸱⸱⸱⸱ \normalfont
\setlength{\parskip}{1.5em}

We love the orange trees. -- \emph{lo mia na lo lipha sathe te lothea}
\setlength{\parskip}{0.3em}

\telcontar ⸱⸱⸱⸱⸱⸱⸱ \normalfont
\setlength{\parskip}{1em}
}
\end{document}
